Aggregation changes two kinds of comparisons: ordering different edges and
ordering the active and inactive windows of one edge. The experiments show
why these effects deserve separate evaluation. Hard aggregation improves
pooled ranking at both duration configurations in Study A, while its
within-edge temporal effect changes sign. Shrinkage achieves pooled gains
without grouping, and residual emission restores comparisons that hard
maps tie. The fixed local scores make these changes attributable to the
emission intervention in the evaluated systems.

The two-state factorisation provides a useful account of relative benefit.
The product of state-assignment discrimination and graph-contrast concordance
must be interpreted against the temporal ranking already present in raw
scores. Better grouping or clearer graph contrast can therefore coexist with
little relative gain when the local scores already discriminate active windows
well. These ground-truth diagnostics explain an observed configuration;
their empirical factors are jointly estimated.

Windowing and execution rules define further conditions of the comparison.
Duration changes at fixed record length alter switch coverage, and Study C
shows substantial variation across window and stride settings under a common
generator. The fixed gate often rejects useful candidates in the longer
records. Static fallback then removes temporal variation altogether, whereas
raw fallback preserves the local ranking on those rejected cases. Reporting
acceptance and the accepted subset's own paired raw comparison makes this
execution effect visible.

For evaluation practice, the results support three complementary reports:
pooled or target-specific structural ranking, within-edge temporal AUROC,
and boundary detection with explicit zero-score and threshold conventions.
Stable-edge changes and gate acceptance complete the operational picture.
The PELT/refit comparison shows the value of this separation for another
partition/refit structure. Reference labels determine which of these reports
are possible: the ESS industrial benchmark supplies a static expert graph,
supporting an external-data alignment illustration in Supplement Section~S7,
while temporal discrimination and boundary recovery require time-indexed truth.

The empirical scope is two-state linear switching systems and local linear
scorers. The clustering experiments use a supplied state count, and the PELT
adapter estimates a segment count from the same local-score features.
Study C's generator modifications are evaluated jointly, so their individual
effects remain unresolved. The auxiliary GMM responsibilities are nearly
binary and provide limited evidence about diffuse soft aggregation.
These conditions identify useful directions for extending the evaluation
protocol to other scorers, state structures and nonlinear systems.

Structural ranking, within-edge temporal discrimination and boundary detection
answer distinct recovery questions. Shared-score controls and diagnostic
accounting make that distinction operational, explaining how an aggregation
rule can improve one answer while weakening another.
